Las señales eléctricas en un circuito electrónico se generan en forma de ondas
progresivas, las cuales se modelan a partir de combinaciones lineales de funciones
$\sin(\omega t)$ y $\cos(\omega t)$, para algún tiempo $t$ y frecuencia angular $\omega$.
Dos señales puramente sinusoidales, coherentes, con misma amplitud $V$ y distinta
frecuencia angular pueden escribirse como
\begin{align*}
	v_1(t) &= V\sin(\omega_1 t) \\
	v_2(t) &= V\sin(\omega_2 t)
\end{align*}
y su superposición, $v_1(t)+v_2(t)$, puede transformarse utilizando identidades
trigonométricas
\begin{equation}\label{eq:super}
	v(t) = 2V\sin\left( \frac{\omega_1+\omega_2}{2} t \right)
			\cos\left( \frac{\omega_1-\omega_2}{2} t \right).
\end{equation}
Es posible identificar dos frecuencias en cada una de las fucnciones trigonométricas.
El primero corresponde al promedio de las frecuencias originales y es la frecuencia
más alta en la expresión. El segundo corresponde a la diferencia media de las
frecuencias y puede interpretarse como la amplitud instantánea de la superposición.

De forma similar, mediante múltiples superposiciones de ondas puramente sinusoidales,
se puede aproximar de forma arbitraria cualquier curva $f(t)$ definida en un periodo
particular utilizando su serie de Fourier. Las señales más convenientes para aproximar
en circuitos eléctricos son la señal cuadrada $C(t)$ y el diente de sierra $S(t)$, ambas
con serie de Fourier
\begin{align}
	C(t) &= V\frac{4}{\pi}\sum_{n=0}^N \frac{\sin( (2n+1)\omega t  )}{2n+1}\label{eq:C} \\
	S(t) &= V\frac{2}{\pi}\sum_{n=1}^N \frac{\sin(n\omega t)}{n}(-1)^{n-1},\label{eq:S}
\end{align}
hasta un término $N$.

Dado que es posible generar las ondas progresivas mencionadas a través de
señales eléctricas, se puede evidenciar experimentalmente sus formas funcionales y
determinar las frecuencias angulares utilizadas. En ese sentido, en el presente informe
se analizó la superposición de dos ondas sinusoidales y los términos principales de
la onda cuadrada y de diente de sierra, con el fin de comprobar las relaciones matemáticas
anteriormente presentadas.